\section{哈代域}

设 \(\mathfrak{F}\) 是 \(\mathbf{R}\) 上由形如 \([x_0, +\infty[\) 的区间组成的滤子基。回忆一下，我们已定义了等价关系 \(R_{\infty}\) ：它在一个实函数所成的集 \(\mathcal{H}(\mathfrak{F}, \mathbf{R})\) 上（这些函数定义于属于 \(\mathfrak{F}\) 的集上），“存在一个集 \(\mathbf{M} \in \mathfrak{F}\) ，使 \(f(x) = g(x)\) 于 \(\mathbf{M}\) 上成立” (V, p.~211)，并且商集 \(\mathcal{H}(\mathfrak{F}, \mathbf{R}) / R_{\infty}\) 带有含单位元的环结构。

定义 1. 给定 \(\mathfrak{K}\) （ \(\mathcal{H}(\mathfrak{F},\mathbf{R})\) 的一个子集），称 \(\mathfrak{K}/\mathrm{R}_{\infty}\) （ \(\mathfrak{K}\) 在 \(\mathcal{H}(\mathfrak{F},\mathbf{R})/\mathrm{R}_{\infty}\) 中的典范像）为一个哈代域，如果 \(\mathfrak{K}\) 满足下列条件：

\(1^{\circ}\) \(\mathfrak{K} / \mathrm{R}_{\infty}\) 是环 \(\mathcal{H}(\mathfrak{F},\mathbf{R}) / \mathrm{R}_{\infty}\) 的一个子域

\(2^{\circ}\) \(\mathfrak{K}\) 中每个函数都在某个区间 \([a, +\infty[\) （随函数而定）上连续且可导，且其导数关于 \(\mathbb{R}_{\infty}\) 的类属于 \(\mathfrak{K}/\mathbb{R}_{\infty}\) 。

\(\mathfrak{K} / \mathrm{R}_{\infty}\) 是一个域这一假设等价于下列条件：若 \(f\in \mathfrak{K}\) 且 \(g\in \mathfrak{K}\) ，则 \(f + g\) 与 \(fg\) 等于 \(\mathfrak{K}\) 中的函数（在 \(\mathfrak{F}\) 的某个集上）；此外，若 \(f\) 在 \(\mathfrak{F}\) 的某个集上不恒为零，则存在 \(\mathfrak{F}\) 中的一个集 M， \(f\) 在 M 上不为零，且 \(1 / f\) 在 M 上等于 \(\mathfrak{K}\) 中的一个函数；由条件 \(2^{\circ}\) ，总可以假定 M 取得使 \(f\) 在 M 上连续，从而在这个区间上具有常号。

（滥用语言的说法）若 \(\mathfrak{K}\) 使 \(\mathfrak{K} / \mathrm{R}_{\infty}\) 是一个哈代域，我们在下文中就说 \(\mathfrak{K}\) 本身是一个哈代域。

例. 1) 每个哈代域都包含有理常数域（特征 0 的最小域，cf.~Alg., V, §1），可以把它与域 \(\mathbf{Q}\) 等同起来；此外，由于两个常数除非相等，否则不能模 \(R_{\infty}\) 同余， \(\mathbf{Q} / R_{\infty}\) 与 \(\mathbf{Q}\) 恒同。实常数也构成一个哈代域，可以把它与 \(\mathbf{R}\) 等同起来。

\begin{enumerate}
\def\labelenumi{\arabic{enumi})}
\setcounter{enumi}{1}
\tightlist
\item
  哈代域的一个非常重要的例子是实系数有理函数的全体，我们（滥用语言的说法）把它记作 \(\mathbf{R}(x)\) ；若 \(f(x) = p(x)/q(x)\) 是一个不恒为零的实系数有理函数，则它连续、可导且 \(\neq 0\) （在一个区间 \([a, +\infty[\) 上），其中 \(a\) 严格大于多项式 \(p(x)\) 与 \(q(x)\) 的实根中最大者；于是 \(\mathbf{R}(x)/\mathrm{R}_{\infty}\) 中除 0 以外的每个元素都可逆。我们再注意，两个有理函数除非相等，否则不能模 \(R_{\infty}\) 同余，故 \(\mathbf{R}(x)/\mathbf{R}_{\infty}\) 仍可与 \(\mathbf{R}(x)\) 等同。
\end{enumerate}

